\documentclass[12pt,a4paper]{article}

\usepackage[utf8]{inputenc}
\usepackage[T1]{fontenc}
\usepackage[french]{babel}
\usepackage{lmodern}
\usepackage[a4paper,margin=2.5cm]{geometry}
\usepackage{hyperref}
\usepackage{xcolor}
\usepackage{array}
\usepackage{booktabs}
\usepackage{tabularx}
\usepackage{enumitem}
\usepackage{amsmath,amssymb}
\usepackage{tcolorbox}
\usepackage{float}
\usepackage{parskip}
\usepackage{tikz}
\usetikzlibrary{arrows.meta,calc}

\definecolor{accent}{HTML}{2A5A8A}
\definecolor{muted}{HTML}{555555}
\definecolor{bgcode}{HTML}{F4F4F4}
\definecolor{warn}{HTML}{B85C00}
\definecolor{thbg}{HTML}{EEF3F8}
\definecolor{okbg}{HTML}{EEF9EE}
\definecolor{ok}{HTML}{2A7A2A}

\hypersetup{colorlinks=true, linkcolor=accent, urlcolor=accent, citecolor=accent}

\tcbuselibrary{skins,breakable}
\newtcolorbox{formulabox}{
  colback=thbg, colframe=accent, boxrule=0.8pt,
  left=8pt, right=8pt, top=6pt, bottom=6pt
}
\newtcolorbox{attention}{
  colback=okbg, colframe=ok, boxrule=0.8pt,
  left=8pt, right=8pt, top=6pt, bottom=6pt,
  title={\textbf{Point clé}}, fonttitle=\small
}

\title{\textbf{Méthodologie et validation terrain}}
\author{Samir El Khoumri, Aly Hachem-Reda, Sami Chbicheb}
\date{Mai 2026}

\begin{document}

\maketitle
\tableofcontents
\newpage

% ============================================================
\section{Méthodologie et validation terrain}
% ============================================================

% ------------------------------------------------------------
\subsection{Le problème de l'évaluation}
% ------------------------------------------------------------

Évaluer la qualité d'un outil de sélection de tests de régression pose
un problème fondamental : l'outil produit une liste de scénarios à
ré-exécuter, mais rien ne permet a priori de savoir si cette liste est
bonne. Pour cela, il faudrait connaître, pour chaque commit, l'ensemble
des scénarios qui \emph{auraient réellement dû} être ré-exécutés.

Ce référentiel de référence est appelé \textbf{vérité terrain}. Sans
lui, il est impossible de distinguer un scénario correctement sélectionné
d'un scénario sélectionné par erreur, ni de détecter les scénarios
impactés que l'outil aurait manqués. La construction de cette vérité
terrain est donc une étape préalable indispensable à toute évaluation
rigoureuse.

% ------------------------------------------------------------
\subsection{Définition de la vérité terrain}
% ------------------------------------------------------------

\subsubsection{Principe}

La vérité terrain repose sur l'observation suivante : \emph{un scénario
est impacté par un commit si et seulement si son exécution traverse au
moins une des méthodes modifiées par ce commit}. Ce principe est à la
fois intuitif et rigoureux --- il s'appuie sur l'exécution réelle des
tests, et non sur une heuristique ou une approximation.

Formellement, pour un diff donné, la vérité terrain $G$ est définie par :

\begin{formulabox}
\[
  G(\mathit{diff}) = \bigl\{\, s \in \mathcal{S}
  \;\mid\;
  \text{Couverture}(s) \cap \text{MéthodesModifiées}(\mathit{diff})
  \neq \emptyset \,\bigr\}
\]
\end{formulabox}

\noindent où $\mathcal{S}$ est l'ensemble des scénarios du projet,
$\text{Couverture}(s)$ l'ensemble des méthodes effectivement exécutées
lors du passage du scénario $s$, et $\text{MéthodesModifiées}(\mathit{diff})$
l'ensemble des méthodes dont le corps est modifié par le diff.

\subsubsection{Couverture dynamique avec JaCoCo}

Pour mesurer $\text{Couverture}(s)$, nous utilisons \textbf{JaCoCo},
un outil d'instrumentation du bytecode Java qui enregistre, lors de
l'exécution d'un test, l'ensemble exact des méthodes traversées.

Un point d'attention important est l'isolation des scénarios : chaque
scénario est exécuté dans une JVM distincte, indépendamment des autres.
Sans cette isolation, la couverture mesurée pour un scénario pourrait
être contaminée par l'exécution d'un scénario précédent (caches en
mémoire, état de la base de données, variables statiques). L'isolation
garantit que la couverture mesurée est bien imputable au seul scénario
considéré.

% ------------------------------------------------------------
\subsection{Métriques d'évaluation}
% ------------------------------------------------------------

Pour évaluer la qualité de la sélection produite par RTS-LLM, on compare
l'ensemble $S$ des scénarios sélectionnés par l'outil à l'ensemble $G$
de la vérité terrain. Trois situations peuvent se présenter pour chaque
scénario :

\begin{table}[H]
\centering
\begin{tabularx}{\linewidth}{llXl}
\toprule
\textbf{Sigle} & \textbf{Nom} & \textbf{Signification} & \textbf{Définition} \\
\midrule
VP & Vrai Positif
   & Le scénario est sélectionné \textbf{et} réellement impacté.
     C'est le cas idéal : l'outil a pris la bonne décision.
   & $|S \cap G|$ \\[6pt]
FP & Faux Positif
   & Le scénario est sélectionné \textbf{mais pas} réellement impacté.
     Le test sera ré-exécuté inutilement, ce qui représente du temps
     de CI perdu, mais ne met pas en danger la qualité.
   & $|S \setminus G|$ \\[6pt]
FN & Faux Négatif
   & Le scénario est impacté \textbf{mais} l'outil ne l'a pas sélectionné.
     C'est le cas le plus grave : un défaut introduit par le commit
     pourrait ne pas être détecté et atteindre la production.
   & $|G \setminus S|$ \\
\bottomrule
\end{tabularx}
\caption{Définition des cas de figure de la sélection.}
\label{tab:vp-fp-fn}
\end{table}

À partir de ces trois quantités, on calcule trois métriques classiques
de la recherche d'information :

\begin{formulabox}
\[
  \text{Précision} = \frac{VP}{VP + FP}
  \qquad
  \text{Rappel} = \frac{VP}{VP + FN}
  \qquad
  F_1 = \frac{2 \times \text{Précision} \times \text{Rappel}}
             {\text{Précision} + \text{Rappel}}
\]
\end{formulabox}

\paragraph{La précision} mesure la proportion de scénarios sélectionnés
qui sont effectivement impactés. Une précision élevée signifie que l'outil
sélectionne peu de scénarios inutiles : le pipeline de CI reste léger.
Une précision faible indique que l'outil est trop conservateur et
sélectionne beaucoup de scénarios sans lien réel avec les modifications.

\paragraph{Le rappel} mesure la proportion de scénarios impactés que
l'outil a effectivement sélectionnés. Un rappel élevé signifie que peu
de scénarios impactés ont été oubliés. C'est \textbf{l'indicateur
prioritaire} pour un outil RTS : un faux négatif (rappel insuffisant)
signifie qu'un scénario impacté n'est pas ré-exécuté, permettant
potentiellement à un défaut d'atteindre la production sans être détecté.
Un faux positif (précision insuffisante), en revanche, n'entraîne qu'une
exécution superflue dont le coût est borné et prévisible.

\paragraph{Le score F1} est la moyenne harmonique de la précision et du
rappel. Il offre une vision synthétique de la qualité globale de la
sélection, pénalisant les outils qui optimiseraient l'une des deux
métriques au détriment de l'autre (par exemple, un outil qui sélectionne
systématiquement tous les scénarios aurait un rappel parfait de 1 mais
une précision très faible).

\begin{attention}
Dans le contexte du RTS, un rappel de 1 est toujours atteignable
trivialement en sélectionnant l'intégralité des scénarios à chaque
commit --- mais cela annule tout l'intérêt de la sélection. L'enjeu est
d'obtenir un rappel élevé \emph{tout en maintenant une précision
raisonnable}, ce que le score F1 permet de mesurer conjointement.
\end{attention}

% ------------------------------------------------------------
\subsection{Mise en place concrète}
% ------------------------------------------------------------

\subsubsection{Collecte de la couverture par scénario}

La première étape consiste à mesurer, une fois pour toutes, la couverture
de chaque scénario dans la version de référence du projet. Pour cela, un
script exécute chaque scénario Cucumber individuellement, avec l'agent
JaCoCo attaché à la JVM. À l'issue de chaque exécution, le fichier de
couverture produit est archivé sous un nom unique lié au scénario.

On obtient ainsi une \textbf{map de couverture} : pour chaque scénario,
la liste des méthodes Java effectivement traversées lors de son passage.

\subsubsection{Calcul de la vérité terrain pour un commit}

Pour un commit donné, le calcul de $G$ suit trois étapes :

\begin{enumerate}[noitemsep]
  \item \textbf{Extraction des méthodes modifiées.} Le diff entre le
        commit et son parent est analysé pour identifier les méthodes
        Java dont le corps a été modifié (par croisement des hunks du
        diff avec les plages de méthodes issues de l'analyse statique).

  \item \textbf{Intersection avec la map de couverture.} Pour chaque
        scénario, on vérifie si l'ensemble de ses méthodes couvertes
        intersecte l'ensemble des méthodes modifiées. Si oui, le scénario
        appartient à $G$.

  \item \textbf{Calcul des métriques.} La sélection $S$ produite par
        RTS-LLM est comparée à $G$ pour calculer VP, FP, FN, précision,
        rappel et F1.
\end{enumerate}

\subsubsection{Application sur l'historique Git}

Pour obtenir un jeu d'évaluation suffisamment représentatif, ce calcul
est automatisé sur une plage de commits de l'historique Git du projet
cible. Pour chaque commit, la couverture est collectée dans l'état du
code \emph{avant} le commit (état parent), afin que les noms et plages
de méthodes soient cohérents avec ce que le diff référence. Les résultats
sont agrégés dans un fichier CSV contenant, pour chaque commit, les
métriques individuelles. Les moyennes finales sont ensuite calculées
sur l'ensemble des commits pour lesquels la vérité terrain est non vide.

% ------------------------------------------------------------
\subsection{Limites connues de la vérité terrain}
% ------------------------------------------------------------

La vérité terrain telle que définie ici est robuste, mais présente trois
cas limites à prendre en compte lors de l'interprétation des résultats.

\paragraph{Méthodes nouvellement introduites par le diff.}
La couverture est collectée avant le commit, dans l'état du code parent.
Si le commit introduit une méthode entièrement nouvelle, aucun scénario
ne peut l'avoir couverte à ce stade : la vérité terrain ne peut pas
identifier les scénarios qui, après le commit, invoqueront cette méthode.
Ce cas tend à \emph{sous-estimer} $G$ et à pénaliser RTS-LLM s'il a
correctement sélectionné ces scénarios.

\paragraph{Fichiers supprimés par le diff.}
Lorsqu'un commit supprime entièrement un fichier, les méthodes qu'il
contenait ne figurent plus dans l'index de l'analyse statique. Elles ne
sont donc pas incluses dans $\text{MéthodesModifiées}(\mathit{diff})$,
même si des scénarios les couvraient dans la version parente.

\paragraph{Surcharges de méthodes.}
Les identifiants de méthodes utilisés n'incluent pas la signature
(types des paramètres). Si une classe possède plusieurs surcharges d'une
même méthode et que seule l'une est modifiée, la vérité terrain peut
inclure des scénarios qui n'utilisaient en réalité qu'une autre surcharge.
Ce cas tend à \emph{surestimer} légèrement $G$.

\medskip
Ces trois limites sont inhérentes à toute approche de vérité terrain par
couverture dynamique et sont documentées dans la littérature
scientifique~\cite{rothermel1997, gligoric2015}. Elles constituent des
biais connus dont il faut tenir compte lors de l'analyse des résultats,
sans remettre en cause la validité globale de la méthodologie.

\begin{thebibliography}{9}
\bibitem{rothermel1997}
  G. Rothermel et M. J. Harrold,
  \emph{A Safe, Efficient Regression Test Selection Technique},
  ACM Transactions on Software Engineering and Methodology, 1997.

\bibitem{gligoric2015}
  M. Gligoric, L. Eloussi et D. Marinov,
  \emph{Practical Regression Test Selection with Dynamic File Dependencies},
  ISSTA, 2015.
\end{thebibliography}

\end{document}
